\specsection{Bibliography}

\begin{enumerate}
\def\labelenumi{(\Roman{enumi})}
\item
  S. BANACH, Sur le problème de la mesure, Fund. Math., 4 (1923), pp.~7--33.
\item
  E. CARTAN, Le principe de dualité et certaines intégrales multiples de l'espace tangentiel et de l'espace réglé, Bull. Soc. Math. France, 24 (1896), pp.~140--177 (= Œuvres complètes, v. II \(_{1}\) , pp.~265--302).
\item
  P. J. DANIELL, Stieltjes--Volterra products, Congr. Intern. des Math., Strasbourg, 1920, pp.~130--136.
\item
  A. HAAR, Der Massbegriff in der Theorie der kontinuierlichen Gruppen, Ann. of Math., (2), 34 (1933), pp.~147--169 (= Gesammelte Arbeiten, pp.~600--622).
\item
  A. HURWITZ, Ueber die Erzeugung der Invarianten durch Integration, Gött. Nachr., 1897, pp.~71--90 (= Math. Werke, v. II, pp.~546--564).
\item
  A. M. MACBEATH, S. SWIERCZKOWSKI, Limits of lattices in a compactly generated group, Canad. J. Math., 12 (1960), pp.~427--437.
\item
  J. von NEUMANN, a) Zum Haarschen Mass in topologischen Gruppen, Comp. Math., 1 (1934), pp.~106--114 (= Collected Works, v. II, n° 22); b) The uniqueness of Haar's measure, Mat. Sbornik, 1 (43) (1936), pp.~721--734 (= Collected Works, v. IV, n° 6).
\item
  P. TCHEBYCHEF, Sur deux théorèmes relatifs aux probabilités, Acta. Math., 14 (1890), pp.~305--315) (= Œuvres, v. II, pp.~481--491).
\item
  V. VOLTERRA, Leçons sur les fonctions de lignes, Paris (Gauthier-Villars), 1913.
\item
  A. WEIL, L'intégration dans les groupes topologiques et ses applications, Actual. Scient. et Ind., n° 869, Paris, Hermann, 1940 (2 \(^{e}\) éd., ibid., n° 869-1145, Paris, Hermann, 1953).
\item
  H. WEYL, a) Theorie des Darstellung kontinuierlicher halbeinfacher Gruppen durch lineare Transformationen, Math. Zeit., 23 (1925), pp.~271--309, 24 (1926), pp.~328--395 and 789--791 (= Selecta, Basel-Stuttgart (Birkhäuser), 1956, pp.~262--366); b) (with F. PETER) Die Vollständigkeit der primitiven Darstellungen einer geschlossenen kontinuierlichen Gruppe, Math. Ann., 97 (1927), pp.~737--755 (= Selecta, pp.~387--404).
\end{enumerate}

